\scanpage{45}
\mspage{22}
ceci est sans doute faux (cf plus bas \inserted{n\textsuperscript{o} 3}) --- il faudrait \inserted{\unsure{en}} \deleted{donc} \inserted{renforcer} c) ci-dessus en

c') L'hom. $E_{K'}^\circ-E_K$ induit par $E_{K'}-E_K$ est anabélien (où $E_{K'}^\circ$ est le noyau de l'hom. composé
\[
E_{K'}^\circ\longrightarrow\Gamma_{\overline{\mathbb Q}/\mathbb Q}\xrightarrow{\substack{\chi\ \text{caractère}\\\text{cyclotomique}}}\widehat{\mathbb Z}^{*}).
\]

Mais pour voir que cette condition est \uline{nécessaire} pour que l'hom. soit géométrique, il faudrait \deleted{savoir} vérifier que \inserted{pour} tout sous-groupe ouvert $E'$ d'un $E_K$, le centralisateur dans $E_K$ (non seulement de $E'$ lui-$\unsure{=|\overline m}$, mais même de $E'^\circ$) est réduit à $1$ --- ce qui résulte de la démonstration du th. 1, et du fait* que pour tt sous-groupe ouvert $\Gamma'$ de $\Gamma=\Gamma_{\mathbb Q}$, le centralisateur (non seulement de $\Gamma'$, mais $\unsure{=|\overline m}$) de $\Gamma'^\circ$ dans $\Gamma$ est réduit à $(1)$.
\margin{* à vérifier!}

Donc, la conjecture initiale ainsi corrigée donne la

\uline{Conséquence} (conjecturale) Pour toute \deleted{groupe de} section de $E_K$ sur \deleted{$\Gamma_{\mathbb Q}$, \unsure{soit}} \inserted{un sous-groupe ouvert $\Gamma'$ de $\Gamma_{\mathbb Q}$}, de sorte que $\Gamma'$ opère (effectivement) sur $\pi_K$, on a \inserted{(si $K$ pas alg. sur $\mathbb Q$ i.e. $\pi_K\ne(1)$)} $\pi_K^{\Gamma'^\circ}\ne(1)$.
\margin{44}
